\reading{MR-10}{Financial Correlation Modeling --- Bottom-Up Approaches}{defer}{2}
% ==========================================================

\begin{mrkeybox}[title={Learning objectives in this reading}]
\small
\begin{enumerate}[label=\textbf{\alph*.},leftmargin=1.7em]
\item Explain the purpose of copula functions and how they are applied in finance.
\item Summarize the process of finding the default time of an asset correlated to all other
      assets in a portfolio using the Gaussian copula.
\item Describe the Gaussian copula and explain how to use it to derive the joint probability of
      default of two assets.
\end{enumerate}
\end{mrkeybox}

\begin{mrtrapbox}[title={Note on ordering}]
The source notes present the Gaussian copula and joint default probability as 9.b and the
correlated default time as 9.c. The 2026 spine reverses them: \textbf{b} is the default time
process, \textbf{c} is the Gaussian copula and joint default probability. The content below
follows the spine.
\end{mrtrapbox}

% ---------------------------------------------------------
\lo{MR-10 a}{Explain the purpose of copula functions and how they are applied in finance.}

\begin{mrdefbox}[title={What a copula is}]
A copula is a mathematical concept that allows us to combine multiple probability distributions
into a single distribution. It takes individual distributions and joins them together to create
a new distribution that describes their relationship. Copulas in finance model joint probability
distributions \term{while preserving the individual distributions}. They map unknown
distributions to a known one, like the uniform or the standard normal distribution.
\end{mrdefbox}

\begin{mrkeybox}[title={Correlation copulas}]
A correlation copula is used to define the correlation between variables or assets. It maps the
individual distributions of the variables to a standard normal distribution, allowing us to
\textit{indirectly} establish a correlation relationship between them.
\end{mrkeybox}

\subsection{Purpose}
The purpose of copula functions is to simplify statistical problems by providing a way to
connect multiple univariate distributions into a single multivariate distribution. They offer a
convenient method for capturing and modeling the dependence or correlation structure between
variables or assets.

\begin{center}
\begin{tikzpicture}[every node/.style={font=\scriptsize\sffamily}]
% --- two arbitrary marginals ---
\begin{scope}[shift={(0,0)}]
  \draw[black!60, line width=0.6pt] (0,0) -- (3.4,0);
  \draw[FRMNavy, line width=1.0pt] plot[smooth, domain=0.1:3.3, samples=60]
    (\x, {1.15*exp(-((\x-1.1)^2)/0.75) + 0.75*exp(-((\x-2.5)^2)/0.30)});
  \node[anchor=south, text=black!70] at (1.7,1.55) {Distribution of $X$};
\end{scope}
\begin{scope}[shift={(4.6,0)}]
  \draw[black!60, line width=0.6pt] (0,0) -- (3.4,0);
  \draw[FRMGreen, line width=1.0pt] plot[smooth, domain=0.1:3.3, samples=60]
    (\x, {0.55*exp(-((\x-0.8)^2)/0.25) + 1.25*exp(-((\x-2.2)^2)/0.85)});
  \node[anchor=south, text=black!70] at (1.7,1.55) {Distribution of $Y$};
\end{scope}
% --- arrows down ---
\draw[FRMGold, -{Stealth[length=5pt]}, line width=1pt] (1.7,-0.25) -- (1.7,-1.15);
\draw[FRMGold, -{Stealth[length=5pt]}, line width=1pt] (6.3,-0.25) -- (6.3,-1.15);
% --- two standard normals ---
\begin{scope}[shift={(0,-2.85)}]
  \draw[black!60, line width=0.6pt] (0,0) -- (3.4,0);
  \draw[FRMNavy, line width=1.0pt] plot[smooth, domain=0.1:3.3, samples=70]
    (\x, {1.35*exp(-((\x-1.7)^2)/0.55)});
  \node[anchor=north, text=black!70] at (1.7,-0.12) {standard normal};
\end{scope}
\begin{scope}[shift={(4.6,-2.85)}]
  \draw[black!60, line width=0.6pt] (0,0) -- (3.4,0);
  \draw[FRMGreen, line width=1.0pt] plot[smooth, domain=0.1:3.3, samples=70]
    (\x, {1.35*exp(-((\x-1.7)^2)/0.55)});
  \node[anchor=north, text=black!70] at (6.3-4.6,-0.12) {standard normal};
\end{scope}
% --- correlation defined here ---
\draw[FRMRed, {Stealth[length=4.5pt]}-{Stealth[length=4.5pt]}, line width=0.9pt]
  (3.5,-2.15) -- (4.5,-2.15);
\node[anchor=west, text=FRMRed, align=left, text width=3.9cm] at (8.25,-2.15)
  {the correlation $\rho$ is defined \textbf{here}, on the normal scale --- never on the
   original distributions};
\draw[FRMRed, -{Stealth[length=4pt]}, line width=0.7pt] (8.15,-2.15) -- (4.7,-2.15);
\end{tikzpicture}
\end{center}
\figcap{Both marginals can be any shape at all. The copula maps each one onto a standard normal
without changing it, and the correlation is then imposed between the two normals. That is the
key property: the original marginal distributions are preserved while a correlation is defined
between them.}

\subsection{How they are applied in finance}
\begin{enumerate}[label=\textbf{\alph*)}]
\item Copulas gained popularity in the financial industry around the year 2000 because they
      aimed to use simple methods to solve complex problems. For example, a copula was assumed
      to need only a single, multidimensional function to correlate assets for a collateralized
      debt obligation structure with 100$+$ assets.
\item They gained traction for simplifying complex dependencies in financial products like CDOs.
\item However, they \term{fell out of favour} during the 2007--2009 financial crisis due to
      limitations in capturing extreme events and systemic risks. Despite this, they are still
      used with caution post-crisis.
\end{enumerate}

\subsection{The copula family}

\begin{center}
\begin{tikzpicture}[every node/.style={font=\scriptsize\sffamily},
  hdrbox/.style={rectangle, rounded corners=2.5pt, draw=FRMNavy!55, fill=PaleNavy,
    line width=0.8pt, text width=3.1cm, align=center, inner sep=4pt, minimum height=0.62cm},
  famb/.style={rectangle, rounded corners=2.5pt, draw=black!40, fill=FRMSoft,
    line width=0.6pt, text width=2.5cm, align=center, inner sep=4pt, minimum height=1.05cm},
  leaf/.style={rectangle, rounded corners=2pt, draw=FRMGold!70, fill=PaleGold,
    line width=0.6pt, text width=1.55cm, align=center, inner sep=3pt, minimum height=0.5cm}]

  \node[hdrbox, text width=6.0cm] (one) at (1.45,0) {\textbf{One-factor copulas}};
  \node[hdrbox, text width=8.6cm] (two) at (9.25,0) {\textbf{Two-factor copulas}};

  \node[famb] (gauss) at (0.05,-1.25) {\textbf{Gaussian}\\ $\rho$ = correlation};
  \node[famb] (arch) at (2.85,-1.25) {\textbf{Archimedean}\\ $\varphi_a(t)$ = generator};

  \node[famb] (stud) at (6.35,-1.25) {\textbf{Student's $t$}\\ $\rho$ = correlation\\ $\nu$ = d.o.f.};
  \node[famb] (frec) at (9.25,-1.25) {\textbf{Fr\'echet}\\ $p$, $q$ = linear combination};
  \node[famb] (mo) at (12.15,-1.25) {\textbf{Marshall-Olkin}\\ $m$, $n$ = weighting factors};

  \node[leaf] (gum) at (3.55,-2.45) {Gumbel};
  \node[leaf] (cla) at (3.55,-3.13) {Clayton};
  \node[leaf] (fra) at (3.55,-3.81) {Frank};

  \draw[black!45, line width=0.6pt] (arch.south) -- (2.45,-2.45) -- (2.45,-3.81);
  \draw[black!45, -{Stealth[length=3.5pt]}, line width=0.6pt] (2.45,-2.45) -- (gum.west);
  \draw[black!45, -{Stealth[length=3.5pt]}, line width=0.6pt] (2.45,-3.13) -- (cla.west);
  \draw[black!45, -{Stealth[length=3.5pt]}, line width=0.6pt] (2.45,-3.81) -- (fra.west);
\end{tikzpicture}
\end{center}
\figcap{The popular copula functions in finance, grouped by how many parameters drive the
dependence. Gaussian is the one the rest of this reading uses; the three Archimedean variants
share a generator function and differ in which tail they concentrate dependence in.}

% ---------------------------------------------------------
\lo{MR-10 b}{Summarize the process of finding the default time of an asset correlated to all other assets in a portfolio using the Gaussian copula.}

The example in the next objective considers only \textit{two} assets. Here we find the default
time for an asset that is correlated to the default times of \textit{all} other assets in a
portfolio.

\begin{mrkeybox}[title={Correlated default time --- the procedure}]
\begin{enumerate}[label=\textbf{\alph*)}]
\item When a Gaussian copula is used to derive the default time relationship for more than two
      assets, a \term{Cholesky decomposition} is used to derive a sample $M_n(\bullet)$ from a
      multivariate copula $M_n(\bullet) \in [0,1]$. The default correlations of the sample are
      determined by the default correlation matrix $\rho_M$ for the $n$-variate standard normal
      distribution $M_n$.
\item Random samples are drawn from that $n$-variate standard normal distribution sample to
      estimate expected default times using the Gaussian copula.
\item Repeat the process of drawing random samples to estimate default times --- for example,
      100,000 times.
\end{enumerate}
\end{mrkeybox}

\begin{mrfmlbox}
\[ M_n(\bullet) \;=\; Q_i(\tau_i) \]
\mrvk{$M_n(\bullet)$ = the value drawn from the multivariate standard normal sample, on
$[0,1]$;\quad $Q_i(\cdot)$ = the cumulative default probability function for asset $i$;\quad
$\tau_i$ = the default time of asset $i$, which is what we are solving for. Read it as: set the
sampled value equal to the cumulative default probability, then invert to get the year.}{}
\end{mrfmlbox}

\begin{center}
\begin{tikzpicture}
\begin{axis}[
  width=10.4cm, height=5.2cm,
  xmin=0.6, xmax=10.4, ymin=0, ymax=0.60,
  xlabel={\scriptsize Default time in years, $\tau_i$},
  ylabel={\scriptsize Default probability $Q_i(\tau_i)$},
  xlabel style={font=\scriptsize}, ylabel style={font=\scriptsize},
  tick label style={font=\scriptsize},
  xtick={1,2,3,4,5,6,7,8,9,10},
  ytick={0,0.1,0.2,0.3,0.4,0.5,0.6},
  axis line style={draw=black!55}, clip=false]
\addplot[black!80, line width=1.1pt, smooth] coordinates {
(1,0.0651)(2,0.1416)(3,0.2103)(4,0.2704)(5,0.3231)(6,0.3673)(7,0.4097)
(8,0.4433)(9,0.4717)(10,0.5001)};
\draw[FRMNavy, -{Stealth[length=4.5pt]}, line width=0.9pt]
  (axis cs:0.65,0.24) -- (axis cs:3.48,0.24);
\draw[FRMNavy, -{Stealth[length=4.5pt]}, line width=0.9pt]
  (axis cs:3.48,0.24) -- (axis cs:3.48,0.008);
\node[circle, fill=FRMNavy, inner sep=1.6pt] at (axis cs:3.48,0.24) {};
\node[anchor=south west, font=\scriptsize\sffamily, text=FRMNavy, fill=white, inner sep=1.5pt] at (axis cs:0.7,0.245)
  {sampled $M_n(\bullet) = 0.24$};
\node[anchor=north west, font=\scriptsize\sffamily, text=FRMNavy, fill=white, inner sep=1.5pt] at (axis cs:3.55,0.055)
  {$\Rightarrow \tau_i \approx 3.5$ years};
\end{axis}
\end{tikzpicture}
\end{center}
\figcap{One draw, read in two moves. Come in horizontally at the sampled value, hit the
cumulative default probability curve, then drop to the axis to read the default time. Repeat
100,000 times and the distribution of those default times is the answer.}

% ---------------------------------------------------------
\lo{MR-10 c}{Describe the Gaussian copula and explain how to use it to derive the joint probability of default of two assets.}

\begin{mrdefbox}[title={The Gaussian copula}]
A Gaussian copula maps the marginal distribution of each variable to the \term{standard normal
distribution}, which by definition has a mean of zero and a standard deviation of one. The key
property of a copula correlation model is preserving the original marginal distributions while
defining a correlation between them.
\end{mrdefbox}

\subsection{The general copula}

\begin{mrfmlbox}
\[ C\bigl[G_1(u_1), \ldots, G_n(u_n)\bigr] \;=\;
   F_n\Bigl[ F_1^{-1}\bigl(G_1(u_1)\bigr), \ldots, F_n^{-1}\bigl(G_n(u_n)\bigr);\ \rho_F \Bigr] \]
\mrvk{$G_i(u_i)$ = the marginal distributions;\quad $F_n$ = the joint cumulative
distribution;\quad $F_1^{-1}$ = the inverse function of $F_n$;\quad $\rho_F$ = the correlation
matrix structure.}{}
\end{mrfmlbox}

\subsection{The Gaussian copula, $C_G$}
Defined for an $n$-variate example, where the joint standard multivariate normal distribution is
denoted $M_n$:

\begin{mrfmlbox}
\[ C_G\bigl[G_1(u_1), \ldots, G_n(u_n)\bigr] \;=\;
   M_n\Bigl[ N_1^{-1}\bigl(G_1(u_1)\bigr), \ldots, N_n^{-1}\bigl(G_n(u_n)\bigr);\ \rho_M \Bigr] \]
\mrvk{$N_i^{-1}$ = the inverse of the univariate standard normal CDF, i.e.\ the percentile
function;\quad $M_n$ = the joint standard multivariate normal distribution;\quad $\rho_M$ = the
correlation matrix of that multivariate normal.}{}
\end{mrfmlbox}

\subsection{The Gaussian default time copula, $C_{GD}$}
In finance, the Gaussian copula is a common approach for measuring default risk. The approach
can be transformed to define the Gaussian default time copula:

\begin{mrfmlbox}
\[ C_{GD}\bigl[Q_i(t), \ldots, Q_n(t)\bigr] \;=\;
   M_n\Bigl[ N_1^{-1}\bigl(Q_1(t)\bigr), \ldots, N_n^{-1}\bigl(Q_n(t)\bigr);\ \rho_M \Bigr] \]
\mrvk{$Q_i(t)$ = the \textit{cumulative} default probability of asset $i$ by time $t$. The only
change from $C_G$ is that the marginals are now cumulative default probability curves.}{}
\end{mrfmlbox}

\begin{mrexbox}[title={Example --- applying a Gaussian copula}]
Two companies, B and Caa, with their estimated default probabilities for years 1 to 10.
\end{mrexbox}

\begin{center}
\small
\begin{tabular}{c r r r r}
\arrayrulecolor{FRMRule}\toprule
\rowcolor{FRMNavy} \hdr{Default time $t$} & \hdr{B default prob.} & \hdr{B cumulative $Q_B(t)$}
  & \hdr{Caa default prob.} & \hdr{Caa cumulative $Q_{Caa}(t)$}\\
1 & 6.51\% & 6.51\% & 23.83\% & 23.83\%\\
\rowcolor{FRMSoft} 2 & 7.65\% & 14.16\% & 13.29\% & 37.12\%\\
3 & 6.87\% & 21.03\% & 10.31\% & 47.43\%\\
\rowcolor{FRMSoft} 4 & 6.01\% & 27.04\% & 7.62\% & 55.05\%\\
5 & 5.27\% & 32.31\% & 5.04\% & 60.09\%\\
\rowcolor{FRMSoft} 6 & 4.42\% & 36.73\% & 5.13\% & 65.22\%\\
7 & 4.24\% & 40.97\% & 4.04\% & 69.26\%\\
\rowcolor{FRMSoft} 8 & 3.36\% & 44.33\% & 4.62\% & 73.88\%\\
9 & 2.84\% & 47.17\% & 2.62\% & 76.50\%\\
\rowcolor{FRMSoft} 10 & 2.84\% & 50.01\% & 2.04\% & 78.54\%\\
\bottomrule
\end{tabular}
\end{center}

\begin{center}
\begin{tikzpicture}
\begin{axis}[
  width=10.4cm, height=5.3cm,
  xmin=0.6, xmax=10.4, ymin=0, ymax=0.85,
  xlabel={\scriptsize Year $t$}, ylabel={\scriptsize Cumulative default probability},
  xlabel style={font=\scriptsize}, ylabel style={font=\scriptsize},
  tick label style={font=\scriptsize},
  xtick={1,2,3,4,5,6,7,8,9,10}, ytick={0,0.2,0.4,0.6,0.8},
  legend style={font=\scriptsize\sffamily, draw=FRMRule, fill=white,
    at={(0.5,-0.30)}, anchor=north, legend columns=-1, column sep=10pt},
  axis line style={draw=black!55}]
\addplot[FRMRed, line width=1.1pt, mark=*, mark size=1.2pt] coordinates {
(1,0.2383)(2,0.3712)(3,0.4743)(4,0.5505)(5,0.6009)(6,0.6522)(7,0.6926)
(8,0.7388)(9,0.7650)(10,0.7854)};
\addlegendentry{$Q_{Caa}(t)$}
\addplot[FRMNavy, line width=1.1pt, mark=*, mark size=1.2pt] coordinates {
(1,0.0651)(2,0.1416)(3,0.2103)(4,0.2704)(5,0.3231)(6,0.3673)(7,0.4097)
(8,0.4433)(9,0.4717)(10,0.5001)};
\addlegendentry{$Q_B(t)$}
\end{axis}
\end{tikzpicture}
\end{center}
\figcap{The two marginals before any copula is applied. Caa is already 23.83\% likely to have
defaulted by the end of year 1, a level B does not reach until year 4. Both curves flatten as
the surviving pool gets smaller.}

\subsection{Step 1: map each cumulative probability to a standard normal percentile}

\begin{center}
\small
\begin{tabular}{c r r r r}
\arrayrulecolor{FRMRule}\toprule
\rowcolor{FRMNavy} \hdr{$t$} & \hdr{$Q_B(t)$} & \hdr{$N^{-1}(Q_B(t))$} & \hdr{$Q_{Caa}(t)$}
  & \hdr{$N^{-1}(Q_{Caa}(t))$}\\
1 & 6.51\% & $-1.5133$ & 23.83\% & $-0.7118$\\
\rowcolor{FRMSoft} 2 & 14.16\% & $-1.0732$ & 37.12\% & $-0.3287$\\
3 & 21.03\% & $-0.8054$ & 47.43\% & $-0.0645$\\
\rowcolor{FRMSoft} 4 & 27.04\% & $-0.6116$ & 55.05\% & 0.1269\\
5 & 32.31\% & $-0.4590$ & 60.09\% & 0.2557\\
\rowcolor{FRMSoft} 6 & 36.73\% & $-0.3390$ & 65.22\% & 0.3913\\
7 & 40.97\% & $-0.2283$ & 69.26\% & 0.5032\\
\rowcolor{FRMSoft} 8 & 44.33\% & $-0.1426$ & 73.88\% & 0.6397\\
9 & 47.17\% & $-0.0710$ & 76.50\% & 0.7225\\
\rowcolor{FRMSoft} 10 & 50.01\% & 0.0003 & 78.54\% & 0.7906\\
\bottomrule
\end{tabular}
\end{center}

\begin{center}
\begin{tikzpicture}
\begin{axis}[
  width=10.6cm, height=5.6cm,
  xmin=-2.6, xmax=2.0, ymin=0, ymax=1.0,
  xlabel={\scriptsize Standard normal percentile $N^{-1}(Q)$},
  ylabel={\scriptsize Cumulative probability $Q$},
  xlabel style={font=\scriptsize}, ylabel style={font=\scriptsize},
  tick label style={font=\scriptsize},
  xtick={-2,-1.5,-1,-0.5,0,0.5,1,1.5,2},
  ytick={0,0.2,0.4,0.6,0.8,1.0},
  axis line style={draw=black!55}, clip=false]
\addplot[black!80, line width=1.1pt, smooth, domain=-2.6:2.0, samples=140]
  {1/(1+exp(-1.702*x))};
\draw[FRMNavy, dashed, line width=0.7pt] (axis cs:-2.6,0.0651) -- (axis cs:-1.5133,0.0651);
\draw[FRMNavy, dashed, line width=0.7pt] (axis cs:-1.5133,0.0651) -- (axis cs:-1.5133,0);
\node[circle, fill=FRMNavy, inner sep=1.5pt] at (axis cs:-1.5133,0.0651) {};
\node[anchor=south west, font=\tiny\sffamily, text=FRMNavy] at (axis cs:-2.55,0.075)
  {$Q_B(1) = 6.51\%$};
\node[anchor=north, font=\tiny\sffamily, text=FRMNavy] at (axis cs:-1.5133,-0.015)
  {$-1.5133$};
\draw[FRMRed, dashed, line width=0.7pt] (axis cs:-2.6,0.2383) -- (axis cs:-0.7118,0.2383);
\draw[FRMRed, dashed, line width=0.7pt] (axis cs:-0.7118,0.2383) -- (axis cs:-0.7118,0);
\node[circle, fill=FRMRed, inner sep=1.5pt] at (axis cs:-0.7118,0.2383) {};
\node[anchor=south west, font=\tiny\sffamily, text=FRMRed] at (axis cs:-2.55,0.248)
  {$Q_{Caa}(1) = 23.83\%$};
\node[anchor=north, font=\tiny\sffamily, text=FRMRed] at (axis cs:-0.7118,-0.015)
  {$-0.7118$};
\end{axis}
\end{tikzpicture}
\end{center}
\figcap{The mapping itself, for year 1. Enter on the vertical axis with the cumulative default
probability, meet the standard normal CDF, then drop to the horizontal axis to read the
percentile. Every row of the table above is one pass through this curve.}

\subsection{Step 2: read the joint probability off the bivariate normal}

\begin{mrexbox}[title={Question 1 --- joint default within the next year}]
What is the joint default probability $Q$ of companies B and Caa in the next year, assuming a
one-year Gaussian default correlation of 0.4?
\tcblower
\small
\[ Q\bigl(t_B \le 1 \cap t_{Caa} \le 1\bigr) \;\equiv\;
   M\bigl( x_B \le -1.5133 \cap x_{Caa} \le -0.7118,\ \rho = 0.4 \bigr) \;=\; \mathbf{3.44\%} \]
\end{mrexbox}

\begin{mrexbox}[title={Question 2 --- different years for each company}]
What is the joint probability of company B defaulting by year 3 and company Caa defaulting by
year 5?
\tcblower
\small
\[ Q\bigl(t_B \le 3 \cap t_{Caa} \le 5\bigr) \;\equiv\;
   M\bigl( x_B \le -0.8054 \cap x_{Caa} \le 0.2557,\ \rho = 0.4 \bigr) \;=\; \mathbf{16.93\%} \]
\end{mrexbox}

\begin{mrkeybox}[title={The shape of the calculation}]
Both questions are the same three moves. Look up each company's \textit{cumulative} default
probability at its own horizon, convert each to a standard normal percentile with $N^{-1}$, then
evaluate the \textit{bivariate} standard normal at those two thresholds with the given
correlation. The correlation enters only at the last step.
\end{mrkeybox}

\begin{mrtrapbox}[title={Two things to watch}]
\begin{itemize}
\item The horizons do not have to match. Question 2 uses year 3 for one company and year 5 for
      the other, and the method is unchanged.
\item The probabilities used are \term{cumulative}, so the result is ``both default \textit{by}
      those dates'', not ``both default \textit{in} those particular years''.
\end{itemize}
\end{mrtrapbox}


% ==========================================================
